\documentclass[twocolumn]{aastex631}
\begin{document}
\title{A residual reionization ionizing-photon-budget shortfall for star-forming galaxies at z\~{}8 (crisis systematic corner)}
\author{NebulaMind Lab (autonomous pipeline)}
\affiliation{NebulaMind Lab, \url{https://lab.nebulamind.net}}
\begin{abstract}
Reionization ionizing-photon-budget reconciliation at z\~{}8: star-forming galaxies require f\_esc=0.789 (+0.628/-0.348) to close the budget (Madau-Dickinson SFRD, log xi\_ion=25.5+/-0.15, clumping C=2-5, JWST-SFRD tail), versus indirect-proxy-inferred f\_esc=0.050 (+0.075/-0.030) from LzLCS O32/beta calibrations. Median delta(required-inferred)=+0.716 dex-frac (16-84\%: +0.361 to +1.350); 99\% of the systematic MC shows a shortfall. Conclusion: a genuine SHORTFALL remains, and the sign holds under both O32 and beta calibrations. The result is bounded by the xi\_ion x clumping x proxy-calibration systematic, not by statistics. Generated autonomously from public data (jwst) via the NebulaMind Lab runner: a bounded, reproducible, descriptive study.
\end{abstract}
\section{Introduction}
Recent studies have highlighted a potential crisis in our understanding of the reionization process, with concerns that star-forming galaxies may not be producing enough ionizing photons to account for the observed reionization [Muñoz2024]. This issue is further complicated by uncertainties in the escape fraction of these photons and the clumping factor of the intergalactic medium [Davies2021]. To address this, we revisit the ionizing photon budget using a literature-anchored approach.
\section{Data and method}
Our calculation relies on the cosmic star formation rate density (SFRD) provided by Madau \& Dickinson's analytic fitting function, along with published values for xi\_ion and O32/beta f\_esc proxy calibrations from LzLCS studies [Chisholm+22, Flury+22; Simmonds+24]. We do not utilize any new survey catalog data or observational results from JWST, SDSS, or TNG. Instead, we focus on reconciling the ionizing photon budget using existing literature values and systematic considerations.
\section{Result}
Our analysis reveals a significant shortfall in the ionizing photon budget at z\~{}8. To close this gap, star-forming galaxies would need an escape fraction of f\_esc=0.789 (+0.628/-0.348), assuming the Madau-Dickinson SFRD, log xi\_ion=25.5+/-0.15, and a clumping factor C between 2-5. However, indirect-proxy-inferred values from LzLCS O32/beta calibrations suggest a much lower escape fraction of f\_esc=0.050 (+0.075/-0.030). The median difference between the required and inferred escape fractions is +0.716 dex-frac (16-84\%: +0.361 to +1.350), with 99\% of systematic Monte Carlo simulations showing a shortfall.
\begin{figure}\centering\includegraphics[width=\columnwidth]{result.png}\caption{A residual reionization ionizing-photon-budget shortfall for star-forming galaxies at z\~{}8 (crisis systematic corner)}\end{figure}
\section{Caveats}
It is essential to acknowledge the limitations of our approach. Our calculation relies on automated, single-selection, uncalibrated measurements, which may not fully capture the complexities of the reionization process. The result is sensitive to the choice of xi\_ion, clumping factor, and proxy calibration, highlighting the need for further observational constraints and refined models to better understand the ionizing photon budget during reionization. Additionally, our analysis does not account for potential systematic errors in the underlying literature values or the impact of other sources of ionizing photons, such as active galactic nuclei.
\section*{References}
\footnotesize
[Muoz2024] Muñoz et al. (2024). Reionization after JWST: a photon budget crisis?. bibcode 2024MNRAS.535L..37M \par [Davies2021] Davies et al. (2021). The Predicament of Absorption-dominated Reionization: Increased Demands on Ionizing Sources. bibcode 2021ApJ...918L..35D \par [Park2022] Park et al. (2022). Calibrating excursion set reionization models to approximately conserve ionizing photons. bibcode 2022MNRAS.517..192P \par [Duncan2015] Duncan et al. (2015). Powering reionization: assessing the galaxy ionizing photon budget at z < 10. bibcode 2015MNRAS.451.2030D \par [Madau2017] Madau (2017). Cosmic Reionization after Planck and before JWST: An Analytic Approach. bibcode 2017ApJ...851...50M

\end{document}
